% =====================================================================
%  SNIPPETS PARA TU main.tex
%  No compiles este archivo: son trozos para copiar y pegar.
% =====================================================================


% ---------------------------------------------------------------------
% (1) AÑADIR AL PREÁMBULO — paquetes que necesita el capítulo 1
% ---------------------------------------------------------------------
\usepackage{enumitem}     % listas con etiquetas personalizadas (OE1, OE2...)
\usepackage{csquotes}     % biblatex se queja si no está, con babel es obligatorio


% ---------------------------------------------------------------------
% (2) ARREGLOS DE LOS AVISOS QUE SALEN EN TU main.log
% ---------------------------------------------------------------------

% 2a. "Package biblatex Warning: 'babel' detected but 'csquotes' missing"
%     -> csquotes SIEMPRE antes de biblatex:
%        \usepackage{csquotes}
%        \usepackage[backend=biber,style=numeric,sorting=none]{biblatex}

% 2b. "Package fancyhdr Warning: \headheight is too small (12.0pt)"
\setlength{\headheight}{14.5pt}
\addtolength{\topmargin}{-2.5pt}

% 2c. Decimales con coma en las cifras de siunitx (memoria en español)
\sisetup{output-decimal-marker={,}}

% 2d. "Font shape T1/lmr/bx/sc undefined" -> versalitas en negrita no existen
%     en Latin Modern. O quitas la negrita, o cargas una fuente que sí las
%     tenga. La opción más simple:
%     \usepackage{newtxtext}   % en lugar de lmodern
%     Si te gusta cómo se ve ahora, ignora el aviso: es cosmético.


% ---------------------------------------------------------------------
% (3) DIVIDIR LA MEMORIA EN VARIOS FICHEROS
%     Estructura de carpetas recomendada:
%
%     memoria/
%       main.tex
%       bibliografia.bib
%       portada.tex
%       capitulos/
%         01_introduccion.tex
%         02_estado_del_arte.tex
%         03_datos.tex
%         04_metodologia.tex
%         05_experimentos.tex
%         06_resultados.tex
%         07_conclusiones.tex
%         A_anexos.tex
%       figuras/
%
%     Y el cuerpo de main.tex queda así:
% ---------------------------------------------------------------------

\begin{document}

\input{portada}                       % \input = pegar tal cual, sin salto de página

\frontmatter                          % (si usas book; con report, usa \pagenumbering{roman})
\tableofcontents
\listoffigures
\listoftables

\mainmatter                           % (con report: \pagenumbering{arabic})

% =====================================================================
%  CAPÍTULO 1 — INTRODUCCIÓN
%  TFG · Decodificación de actividad neuronal a texto mediante la
%        adaptación de un modelo de habla preentrenado (OWSM)
% =====================================================================
% NOTA DE TERMINOLOGÍA (leer antes de tocar nada):
% Los datos del Brain-to-Text '25 NO son sEEG. Son registros
% intracorticales obtenidos con cuatro matrices de microelectrodos de
% alta densidad (64 canales cada una, 256 electrodos en total) en la
% corteza motora del habla. Si en algún sitio has escrito sEEG,
% cámbialo. Ver la nota que te he dejado en el chat.
% =====================================================================

\chapter{Introducción}
\label{cap:introduccion}

\section{Contexto}
\label{sec:contexto}

La comunicación verbal es una de las capacidades más básicas de la vida en sociedad y, sin embargo, puede perderse por completo mientras el resto de las funciones cognitivas permanecen intactas. Enfermedades neurodegenerativas como la esclerosis lateral amiotrófica (ELA), los ictus de tronco encefálico o determinadas lesiones medulares altas pueden provocar disartria severa, anartria o incluso el síndrome de cautiverio (\textit{locked-in syndrome}): un estado en el que la persona conserva la conciencia y el lenguaje, pero carece de la vía motora necesaria para articularlo \cite{card2024}. El intento de hablar sigue produciéndose en la corteza motora; lo que falla es el músculo, no el mensaje.

Las tecnologías de apoyo a la comunicación actualmente disponibles, como los teclados controlados por seguimiento ocular, permiten sortear parcialmente esa barrera, pero a un coste elevado en velocidad y esfuerzo: sus tasas de comunicación se sitúan en el orden de unas pocas palabras por minuto, frente a las aproximadamente 150 palabras por minuto del habla conversacional \cite{willett2023}. Esa diferencia de más de un orden de magnitud no es un detalle técnico, sino la distancia que separa el poder deletrear una necesidad del poder mantener una conversación.

Las interfaces cerebro-computador (\textit{Brain-Computer Interfaces}, BCI) abordan el problema desde el otro extremo del canal: en lugar de aprovechar el residuo motor que le quede al usuario, registran directamente la actividad cortical asociada a la intención y la traducen a una acción externa. Una BCI se compone, de forma general, de cuatro etapas: adquisición de la señal neuronal, extracción de características, decodificación mediante un modelo aprendido y realimentación al usuario. Las tecnologías de adquisición se ordenan en un compromiso claro entre invasividad y calidad de señal, que va desde la electroencefalografía (EEG) de superficie, no invasiva pero de baja relación señal-ruido y escasa resolución espacial, hasta los registros intracorticales con matrices de microelectrodos implantadas en el tejido, que resuelven la actividad de poblaciones neuronales locales y son, hoy por hoy, la única modalidad con la que se han alcanzado prestaciones compatibles con una conversación \cite{willett2023,card2024}.

% TODO: si quieres, aquí encaja una figura del pipeline genérico de una BCI
% (adquisición -> características -> decodificador -> salida).

\section{Del control motor al habla: los sistemas \textit{brain-to-text}}
\label{sec:brain-to-text}

Durante la última década, las BCI intracorticales han pasado de controlar cursores y brazos robóticos a decodificar lenguaje. Un primer salto cualitativo se produjo al decodificar la escritura manual imaginada, alcanzando unos 90 caracteres por minuto en un participante con tetraplejia \cite{willett2021}. Poco después, dos trabajos simultáneos demostraron la decodificación del habla intentada: uno mediante electrocorticografía \cite{metzger2023} y otro mediante matrices intracorticales en la corteza premotora ventral, con una tasa de error por palabra (\textit{Word Error Rate}, WER) del \SI{23.8}{\percent} sobre un vocabulario de 125.000 palabras a unas 62 palabras por minuto \cite{willett2023}.

El referente actual, y el que da origen a los datos empleados en este trabajo, es la neuroprótesis descrita por Card \textit{et al.} \cite{card2024}: en un participante con ELA avanzada y disartria severa, el sistema alcanzó un \SI{99.6}{\percent} de acierto con un vocabulario de 50 palabras el primer día de uso y un \SI{90.2}{\percent} sobre un vocabulario de 125.000 palabras el segundo día, sosteniendo posteriormente en torno al \SI{97.5}{\percent} de acierto durante más de ocho meses y permitiendo conversaciones autónomas a unas 32 palabras por minuto a lo largo de más de 248 horas acumuladas de uso. Estos resultados sitúan el problema en una fase distinta: la viabilidad clínica ya no está en discusión, y la pregunta se desplaza hacia la eficiencia, la reproducibilidad y la generalización de los decodificadores.

Resulta especialmente relevante para este trabajo la forma que adopta la arquitectura de estos sistemas. Prácticamente todos ellos comparten la misma estructura: una secuencia de características neuronales muestreada a intervalos regulares se procesa mediante una red recurrente o un \textit{transformer}, se entrena con una pérdida CTC (\textit{Connectionist Temporal Classification}) \cite{graves2006ctc} para emitir una secuencia de fonemas, y esa secuencia se convierte en texto mediante un modelo de lenguaje $n$-grama y, en los sistemas más recientes, un reordenamiento posterior con un modelo de lenguaje de gran tamaño. Es, punto por punto, la arquitectura clásica del reconocimiento automático del habla (\textit{Automatic Speech Recognition}, ASR), con la única diferencia de que la entrada no es un espectrograma sino actividad cortical.

Esa coincidencia estructural plantea una pregunta natural que este trabajo toma como punto de partida. En procesamiento del lenguaje, visión por computador y reconocimiento del habla, el paradigma dominante consiste en preentrenar un modelo de gran capacidad sobre cantidades masivas de datos y adaptarlo después a la tarea concreta con una fracción mínima de datos y cómputo. En decodificación neuronal, en cambio, lo habitual sigue siendo entrenar el decodificador desde cero para cada participante, precisamente en el escenario donde los datos son más caros de obtener: cada hora de entrenamiento exige la colaboración activa de una persona enferma en una sesión clínica supervisada. Si un codificador de habla preentrenado hubiera aprendido algo genuinamente transferible sobre la estructura temporal de la producción del lenguaje, ese conocimiento debería poder reutilizarse aunque la modalidad de entrada cambie por completo.

\section{Planteamiento del trabajo}
\label{sec:planteamiento}

Este Trabajo de Fin de Grado estudia la viabilidad de reutilizar un modelo de reconocimiento del habla preentrenado a gran escala como codificador de señal neuronal para una tarea \textit{brain-to-text}. Concretamente, se adapta OWSM v3.1 (\textit{Open Whisper-style Speech Model}) \cite{peng2024owsmv31}, una alternativa abierta y reproducible a Whisper \cite{radford2023whisper} construida sobre la arquitectura E-Branchformer \cite{kim2023ebranchformer}, para decodificar a texto la actividad intracortical del conjunto de datos público del Brain-to-Text~'25 \cite{b2t25,card2024}, registrada en el participante T15 mediante 256 microelectrodos a lo largo de 45 sesiones y 20 meses.

La adaptación requiere resolver un desajuste de modalidad: el modelo espera un banco de filtros mel derivado de audio y recibe, en su lugar, un vector de 512 características neuronales por ventana temporal. Para salvarlo se sustituye el \textit{frontend} acústico original por un módulo adaptador entrenado desde cero, se antepone una transformación afín específica de cada sesión que absorbe la deriva de la señal entre días de registro, y se entrena el conjunto con un objetivo CTC puro, prescindiendo del decodificador autorregresivo del modelo original. Sobre esa base se estudian, mediante ablaciones controladas, las cuatro decisiones de diseño que determinan el comportamiento del sistema: cuánto aporta realmente el preentrenamiento, a qué granularidad lingüística conviene leer la corteza, cuántos parámetros hace falta adaptar y qué parte del conocimiento lingüístico acaba residiendo en el decodificador y qué parte en el modelo de lenguaje externo.

Conviene precisar desde el principio el carácter del trabajo. No se persigue competir con los sistemas que encabezan el \textit{benchmark} oficial, que recurren a conjuntos de decenas de modelos y a modelos de lenguaje de gran tamaño para la fusión final de hipótesis, y cuyos requisitos de cómputo y memoria exceden con mucho los recursos disponibles aquí. El objetivo es distinto y complementario: producir evidencia experimental limpia y reproducible sobre \textit{qué transfiere y qué no} de un modelo de habla a una modalidad que no es habla, obtenida íntegramente con herramientas de código abierto, datos públicos y una única GPU de gama media.

\section{Motivación}
\label{sec:motivacion}

\subsection{Motivación clínica y social}

La ELA presenta una incidencia aproximada de 2 casos por cada 100.000 habitantes y año en poblaciones europeas, y la mayoría de las personas afectadas desarrolla algún grado de deterioro del habla a lo largo de la enfermedad \cite{longinetti2019}. % TODO: verificar esta cifra y la referencia antes de la entrega.
A esta población se suman los supervivientes de ictus de tronco y otras lesiones neurológicas que cursan con anartria. Para todas ellas, la pérdida del habla no supone únicamente una limitación funcional: supone la pérdida de la autonomía en la relación con los demás,.
% TODO (Kiril): aqui tenia una frase diciendo que recuperar el habla aparece
% sistematicamente entre las prioridades de los pacientes. Es cierto y es un
% buen argumento, pero necesita cita (hay encuestas de prioridades en ELA y en
% lesion medular). La he quitado hasta que la tengas; si la recuperas, citala.

Las neuroprótesis del habla han demostrado ya que esa capacidad puede restituirse con calidad conversacional \cite{card2024}. El cuello de botella se ha desplazado, por tanto, hacia cuestiones de ingeniería que condicionan su despliegue real: cuántos datos de calibración necesita el sistema, cuánto tarda en estar operativo tras el implante, y en qué medida sigue funcionando en una sesión nueva, un día distinto, sin recalibración. Cualquier avance que reduzca la cantidad de datos que hay que pedirle a un paciente enfermo tiene un impacto clínico directo, y el aprendizaje por transferencia es, precisamente, la familia de técnicas diseñada para ello.

\subsection{Motivación científica y técnica}

Desde el punto de vista técnico, el interés del trabajo reside en que plantea una pregunta que rara vez se formula de manera explícita: ¿qué es exactamente lo que aprende un modelo de habla preentrenado? La respuesta habitual, que aprende a reconocer voz, es demasiado gruesa. Un modelo como OWSM contiene, al menos, tres tipos de conocimiento superpuestos: un \textit{frontend} acústico especializado en la señal de audio, un codificador que modela dependencias temporales de largo alcance en secuencias con estructura de lenguaje, y una cabeza de salida sobre un vocabulario de subpalabras que incorpora estadística ortográfica del inglés. Sustituir la entrada por actividad cortical desacopla estos tres componentes y permite medirlos por separado. Que las unidades de salida a nivel de carácter llegasen a superar a las subpalabras preentrenadas, por ejemplo, no sería un contratiempo experimental, sino una medida informativa sobre dónde reside y dónde no reside el conocimiento transferible.

Existe además un argumento neurocientífico que estructura buena parte de los experimentos. La corteza motora del habla codifica articulación, no ortografía. La unidad de salida a priori correcta es, por tanto, el fonema, y de hecho es la que utilizan tanto el sistema de referencia de este conjunto de datos como la práctica totalidad de las entradas del \textit{benchmark}, que seleccionan modelo por tasa de error de fonema. Comparar sistemáticamente distintas granularidades de salida sobre un mismo eje de evaluación permite contrastar esa hipótesis con datos propios en lugar de asumirla.

\subsection{Motivación de accesibilidad y reproducibilidad}

Los sistemas que encabezan actualmente el \textit{benchmark} combinan decenas de modelos entrenados con semillas distintas y emplean modelos de lenguaje de gran tamaño para la fusión de sus salidas, con unos requisitos de cómputo y de memoria que se detallan en el \cref{cap:estado-del-arte}. Este trabajo se desarrolla íntegramente sobre datos públicos, bibliotecas de código abierto y una única GPU accesible mediante un servicio en la nube de bajo coste, y documenta de forma explícita el presupuesto de cómputo de cada experimento. Mostrar qué parte del problema sigue siendo abordable en esas condiciones tiene valor por sí mismo, tanto para grupos de investigación con recursos limitados como para el propio criterio de reproducibilidad de los resultados publicados.

\subsection{Motivación personal}

% TODO (Kiril): reescribe este párrafo entero con tus palabras. Te dejo un
% andamiaje para que no partas de cero, pero esto tiene que sonar a ti.
El interés personal por este trabajo surge de la confluencia de dos ámbitos que hasta ahora había abordado por separado a lo largo del grado: el procesamiento de señal y el aprendizaje profundo aplicado a secuencias, por un lado, y la aplicación de ambos a un problema con una finalidad médica tangible, por otro. Enfrentarse a un problema en el que el modelo no funciona hasta que se entienden simultáneamente la física de la señal, la arquitectura del modelo y la naturaleza lingüística de la salida ha resultado ser, además, la mejor forma posible de cerrar la formación del grado.

\section{Objetivos}
\label{sec:objetivos}

\subsection{Objetivo general}

El objetivo general de este trabajo es \textbf{evaluar la viabilidad de reutilizar un modelo de reconocimiento del habla preentrenado a gran escala (OWSM v3.1) como codificador de actividad neuronal intracortical en una tarea de decodificación \textit{brain-to-text}}, cuantificando de forma experimental qué componentes de dicho preentrenamiento transfieren a una modalidad no acústica, bajo qué configuración de adaptación lo hacen, y cómo se reparte el conocimiento lingüístico entre el decodificador neuronal y el modelo de lenguaje externo.

\subsection{Preguntas de investigación}

El objetivo general se concreta en tres preguntas de investigación que vertebran el diseño experimental y la discusión de resultados:

\begin{description}
  \item[PI1. Transferencia.] ¿Qué parte de un modelo de habla preentrenado transfiere a una modalidad que no es habla? Es decir, ¿aporta el preentrenamiento de OWSM una ventaja medible frente a la misma arquitectura inicializada aleatoriamente, y cuántos parámetros del codificador es necesario adaptar para materializarla?
  \item[PI2. Granularidad.] ¿A qué granularidad lingüística puede leerse la corteza motora del habla? ¿Cómo se comportan las unidades de salida de carácter, fonema, subpalabra y palabra cuando se comparan sobre un mismo eje de evaluación, y qué papel juega en ello el número de ejemplos de entrenamiento disponibles por clase?
  \item[PI3. Localización del conocimiento lingüístico.] ¿Dónde reside el conocimiento lingüístico del sistema completo, en el codificador neuronal o en el modelo de lenguaje externo? ¿Es el modelo de lenguaje un mero reordenador de hipótesis acústicas o debe guiar activamente la búsqueda durante la decodificación?
\end{description}

\subsection{Objetivos específicos}

Para dar respuesta a lo anterior se plantean los siguientes objetivos específicos, cada uno asociado a un bloque experimental concreto:

\begin{enumerate}[label=\textbf{OE\arabic*.}, leftmargin=*, itemsep=0.4em]

  \item \textbf{Construir una cadena de procesado y entrenamiento reproducible} que permita ingerir los registros neuronales del Brain-to-Text~'25 en formato HDF5 y entrenar sobre ellos un modelo OWSM mediante ESPnet-EZ \cite{someki2024espnetez,watanabe2018espnet}, incluyendo la verificación de integridad de los datos, la parametrización completa de cada experimento mediante identificadores trazables y el registro sistemático de resultados y curvas de entrenamiento.

  \item \textbf{Diseñar el módulo de adaptación de modalidad} que sustituya al \textit{frontend} acústico del modelo, de forma que las 512 características neuronales por ventana temporal se proyecten al espacio de entrada esperado por el codificador, e incorporar una capa específica de sesión, en forma de transformación afín aprendida por cada día de registro, que compense la deriva de la señal entre sesiones.

  \item \textbf{Cuantificar la aportación real del preentrenamiento} (PI1) mediante un control científico explícito que compare, con presupuesto de cómputo idéntico, tres condiciones: codificador preentrenado y ajustado, codificador preentrenado y congelado, y codificador con inicialización aleatoria y ajustado.

  \item \textbf{Determinar la granularidad de salida más adecuada} (PI2) comparando objetivos CTC de carácter, fonema, subpalabra específica del corpus, palabra de vocabulario cerrado y subpalabra original del modelo preentrenado, evaluando todas las variantes sobre la tasa de error por palabra del texto finalmente decodificado, que constituye el único eje de comparación común entre unidades distintas.

  \item \textbf{Caracterizar el compromiso entre capacidad y coste del módulo adaptador}, comparando arquitecturas de complejidad creciente bajo una receta de entrenamiento idéntica y reportando, junto al rendimiento, el número de parámetros y el tiempo de entrenamiento por época, de modo que las comparaciones queden libres de las confusiones experimentales presentes en las primeras iteraciones del trabajo.

  \item \textbf{Analizar la relación entre parámetros entrenables y rendimiento} (PI1) mediante estrategias de adaptación de coste creciente, desde la adaptación de bajo rango (LoRA) \cite{hu2022lora} sobre un codificador congelado hasta el descongelado progresivo de sus capas superiores, obteniendo una curva que permita identificar el punto de operación más eficiente.

  \item \textbf{Evaluar la generalización a sesiones no observadas}, escenario que reproduce la situación real de despliegue clínico, midiendo la degradación del sistema sobre sesiones excluidas del entrenamiento y cuantificando el compromiso que introduce la capa específica de sesión, así como el coste de una adaptación mínima a una sesión nueva.

  \item \textbf{Determinar el papel del modelo de lenguaje externo} (PI3) integrando un decodificador CTC por búsqueda en haz con un modelo de lenguaje $n$-grama \cite{heafield2011kenlm}, explorando el espacio de sus hiperparámetros de fusión, comparando distintos órdenes del modelo y contrastando su rendimiento con el límite superior que impone la lista de hipótesis acústicas (análisis \textit{oracle}), lo que permite discriminar si el modelo de lenguaje reordena o guía la búsqueda.

  \item \textbf{Consolidar la configuración final y estimar su variabilidad}, entrenando el mejor sistema resultante con varias semillas de inicialización para reportar los resultados con una medida de dispersión que permita distinguir las diferencias reales entre configuraciones del ruido experimental, y situar dichos resultados en el contexto del sistema de referencia oficial del conjunto de datos.

\end{enumerate}

\subsection{Alcance y limitaciones}
\label{sec:alcance}

Para delimitar con precisión qué se puede y qué no se puede concluir de este trabajo, se hacen explícitas desde el inicio las siguientes restricciones:

\begin{itemize}
  \item \textbf{Un único participante.} Todos los experimentos se realizan sobre los datos del participante T15. Las conclusiones sobre transferencia y granularidad son, por tanto, específicas de este sujeto, de la localización de sus implantes y del inglés como lengua objetivo, y no pueden extrapolarse sin más a otros participantes o idiomas.
  \item \textbf{Decodificación fuera de línea.} El sistema se evalúa en diferido sobre grabaciones ya adquiridas. No se aborda la ejecución en tiempo real ni la latencia percibida por el usuario, ambas condiciones necesarias para un uso conversacional.
  \item \textbf{Vocabulario cerrado en una de las configuraciones.} El objetivo de salida a nivel de palabra, que resulta ser el de mejor rendimiento, sólo puede emitir palabras vistas en entrenamiento. Sus cifras no son extrapolables a un escenario de vocabulario abierto.
  \item \textbf{Presupuesto de cómputo acotado.} El trabajo se ejecuta sobre una única GPU de gama media en un entorno en la nube. Esto excluye por construcción las técnicas de conjunto de decenas de modelos y el uso de modelos de lenguaje de gran tamaño para la fusión de hipótesis, que son las que sostienen los primeros puestos del \textit{benchmark}. Los resultados absolutos deben leerse teniendo esto presente.
  \item \textbf{Modelo de lenguaje propio.} El modelo de lenguaje $n$-grama oficial de la competición requiere una cantidad de memoria principal que excede la disponible en el entorno de trabajo, por lo que se entrena un modelo propio de orden reducido sobre las transcripciones del corpus. Esta restricción condiciona los valores absolutos de tasa de error por palabra, aunque no invalida las comparaciones relativas, que se realizan siempre con el mismo modelo de lenguaje.
  \item \textbf{Ablaciones con presupuesto igualado, no hasta convergencia.} Las comparaciones entre configuraciones se realizan asignando a todas las variantes el mismo número de épocas y desactivando la parada temprana, de modo que la comparación sea justa. Los valores reportados en las ablaciones no corresponden, por tanto, al rendimiento máximo alcanzable por cada configuración, sino a su rendimiento bajo un presupuesto común.
\end{itemize}

\section{Estructura de la memoria}
\label{sec:estructura}

% TODO: ajusta esta lista a los títulos definitivos de tus capítulos.
El resto de la memoria se organiza como sigue. El \cref{cap:estado-del-arte} revisa el estado del arte en neuroprótesis del habla y en la adaptación de modelos de habla preentrenados, y sitúa la aportación de este trabajo respecto a ambas líneas. El \cref{cap:datos} justifica la elección del conjunto de datos frente a las alternativas evaluadas, lo describe en detalle y analiza los retos que impone. El \cref{cap:marco-teorico} desarrolla los fundamentos de los modelos y algoritmos empleados: el codificador atencional, el objetivo CTC, la adaptación de rango bajo, la decodificación con modelo de lenguaje y las métricas. El \cref{cap:arquitectura} describe el sistema propuesto, detallando qué se conserva del modelo preentrenado, qué se sustituye y qué se añade. El \cref{cap:metodologia} expone el protocolo experimental, de evaluación y de comparación, junto con el diseño de los bloques de experimentos. El \cref{cap:resultados} recoge y discute los resultados obtenidos. El \cref{cap:limitaciones} reúne las limitaciones del estudio y las líneas de continuación que se consideran más prometedoras. Finalmente, el \cref{cap:conclusiones} resume las conclusiones del trabajo.
   % \include = capítulo, empieza en página nueva
% =====================================================================
%  CAPÍTULO 2 — ESTADO DEL ARTE
% =====================================================================
% Requiere en el preámbulo: booktabs, enumitem, cleveref, siunitx
% =====================================================================

\chapter{Estado del arte}
\label{cap:estado-del-arte}

Este capítulo revisa el contexto científico y técnico en el que se inscribe el trabajo. Se organiza en tres bloques. El primero (\cref{sec:sota-bci,sec:sota-neuroprotesis,sec:sota-anatomia,sec:sota-benchmarks}) recorre la evolución de las neuroprótesis del habla hasta los sistemas actuales y disecciona la arquitectura que todos ellos comparten. El segundo (\cref{sec:sota-habla,sec:sota-peft}) resume el estado del arte en modelos de habla preentrenados y en las técnicas de adaptación eficiente que hacen viable su reutilización. El tercero (\cref{sec:sota-transferencia,sec:sota-posicionamiento}) revisa los trabajos que han intentado explícitamente cruzar ambos campos y sitúa la aportación de esta memoria respecto a ellos.

\section{Modalidades de registro en interfaces cerebro-computador}
\label{sec:sota-bci}

La elección de la técnica de registro condiciona por completo el tipo de información que puede extraerse y, en consecuencia, la clase de decodificador que tiene sentido construir. Las modalidades empleadas en BCI del habla se ordenan en un compromiso entre invasividad y calidad de señal:

\begin{description}
  \item[Electroencefalografía (EEG).] No invasiva, de bajo coste y ampliamente disponible. Registra la actividad sumada de grandes poblaciones neuronales a través del cráneo, lo que impone una resolución espacial de centímetros y una relación señal-ruido baja. Es adecuada para paradigmas de selección discreta, pero los intentos de decodificar habla continua a partir de EEG han producido resultados cuya validez ha sido cuestionada de forma reiterada, en buena medida por problemas de partición de datos y por métricas que premian la similitud semántica en lugar de la recuperación literal del estímulo.
  \item[Estereoelectroencefalografía (sEEG).] Electrodos de profundidad implantados con fines diagnósticos, habitualmente en pacientes con epilepsia refractaria. Ofrecen buena resolución temporal y acceso a estructuras profundas, pero su colocación responde a criterios clínicos y no a la conveniencia del decodificador, y los registros suelen ser de corta duración.
  \item[Electrocorticografía (ECoG).] Rejillas de electrodos colocadas sobre la superficie cortical, sin penetrar en el tejido. Combinan una cobertura espacial amplia con una estabilidad crónica notable, y sobre ellas se construyeron los primeros sistemas de habla en tiempo real \cite{moses2021,metzger2023}.
  \item[Matrices intracorticales de microelectrodos.] Conjuntos de electrodos penetrantes (típicamente matrices tipo Utah) que registran la actividad de poblaciones neuronales locales con resolución de unidad múltiple. Es la modalidad de mayor riesgo quirúrgico y, a la vez, la única con la que se han alcanzado prestaciones compatibles con la conversación \cite{willett2023,card2024}. Es también la modalidad de los datos empleados en este trabajo.
\end{description}

De las características extraídas de estos registros, dos se han impuesto como estándar en las matrices intracorticales: el recuento de cruces de umbral (\textit{threshold crossings}), que cuenta los eventos en que la señal filtrada cae por debajo de un umbral proporcional al valor eficaz del ruido, y la potencia en la banda de espigas (\textit{spike band power}). Ambas se agregan en ventanas de \SI{20}{\milli\second}, lo que produce una secuencia a \SI{50}{\hertz} \cite{b2t25}.

\section{Evolución de las neuroprótesis del habla}
\label{sec:sota-neuroprotesis}

Los primeros sistemas de comunicación basados en BCI operaban por selección: el usuario elegía letras o comandos de un conjunto discreto. El salto hacia la producción continua de lenguaje llegó por dos vías simultáneas.

La primera fue la decodificación de la escritura manual imaginada, que alcanzó unos 90 caracteres por minuto con una precisión superior al \SI{99}{\percent} empleando corrección automática \cite{willett2021}. Aunque el gesto decodificado no es el habla, ese trabajo introdujo dos elementos que siguen siendo estándar: el entrenamiento con pérdida CTC sobre secuencias sin alineamiento explícito, y las capas de entrada específicas de cada día de registro para compensar la deriva de la señal.

La segunda vía atacó directamente el habla intentada. Moses \textit{et al.} \cite{moses2021} demostraron la decodificación en tiempo real en una persona con anartria mediante ECoG, con un vocabulario cerrado de 50 palabras. Dos años después, dos trabajos publicados simultáneamente ampliaron el problema a vocabulario abierto: Metzger \textit{et al.} \cite{metzger2023} con ECoG de alta densidad, incorporando síntesis de voz y control de un avatar, y Willett \textit{et al.} \cite{willett2023} con matrices intracorticales en la corteza premotora ventral, alcanzando un \SI{23.8}{\percent} de tasa de error por palabra sobre un vocabulario de 125.000 palabras a unas 62 palabras por minuto.

El referente actual es el sistema descrito por Card \textit{et al.} \cite{card2024}, del que proceden los datos de este trabajo. Sus resultados marcan un cambio de fase respecto a los anteriores: \SI{99.6}{\percent} de acierto con vocabulario de 50 palabras el primer día de uso, \SI{90.2}{\percent} sobre 125.000 palabras el segundo día tras \num{1.4} horas adicionales de datos, y en torno al \SI{97.5}{\percent} de acierto sostenido durante más de ocho meses, con más de 248 horas acumuladas de conversación autónoma a unas 32 palabras por minuto. La \cref{tab:sota-sistemas} resume esta progresión.

\begin{table}[htbp]
  \centering
  \caption{Progresión de los sistemas de decodificación de habla de vocabulario abierto. WER: tasa de error por palabra. PPM: palabras por minuto.}
  \label{tab:sota-sistemas}
  \small
  \begin{tabular}{llllll}
    \toprule
    Trabajo & Registro & Participante & Vocabulario & WER & PPM \\
    \midrule
    Moses \textit{et al.} \cite{moses2021}    & ECoG          & Anartria        & 50 palabras   & ---            & $\sim$15 \\
    Willett \textit{et al.} \cite{willett2021}& Intracortical & Tetraplejia     & Escritura     & $<$\,1\,\%$^{*}$ & 90 cpm \\
    Metzger \textit{et al.} \cite{metzger2023}& ECoG          & Ictus de tronco & 1.024 palabras& ---            & $\sim$78 \\
    Willett \textit{et al.} \cite{willett2023}& Intracortical & T12 (ELA)       & 125.000       & 23,8\,\%       & 62 \\
    Card \textit{et al.} \cite{card2024}      & Intracortical & T15 (ELA)       & 125.000       & $\sim$2,5\,\%  & 32 \\
    \bottomrule
  \end{tabular}

  \vspace{0.4em}
  {\footnotesize $^{*}$ Tasa de error por carácter con corrección automática, no WER. CPM: caracteres por minuto.}
  % TODO: verifica las cifras de Moses y Metzger antes de la entrega; las he
  % dejado aproximadas a propósito. La de Willett 2021 es tasa de error por
  % carácter con autocorrección, no WER: ojo al pie de tabla.
\end{table}

\section{Anatomía de un decodificador \textit{brain-to-text}}
\label{sec:sota-anatomia}

Pese a la diversidad de grupos y modalidades, los sistemas anteriores comparten una arquitectura común que conviene desglosar, porque es exactamente la arquitectura de un sistema de reconocimiento automático del habla y es esa coincidencia la que motiva este trabajo.

\subsection{Características de entrada y no estacionariedad}

La entrada es una secuencia de vectores de características neuronales a \SI{50}{\hertz}. Un problema transversal, y probablemente el más limitante en la práctica, es la \textbf{no estacionariedad}: la relación entre actividad registrada e intención cambia entre sesiones por micromovimientos del implante, gliosis, variaciones en el estado del participante y degradación progresiva del electrodo. Un decodificador entrenado con datos de una sesión se degrada de forma apreciable en la siguiente. La solución estándar, introducida en \cite{willett2021} y presente en prácticamente todos los sistemas posteriores, consiste en aprender una transformación de entrada específica de cada sesión (típicamente afín) que proyecta los datos de cada día a un espacio común, dejando que el resto del modelo sea compartido. Los sistemas más recientes refinan esta idea con proyecciones jerárquicas de bajo rango a nivel de mes y de día \cite{brainwhisperer}.

\subsection{El codificador neuronal}

El codificador transforma la secuencia de características en una secuencia de distribuciones sobre unidades lingüísticas. El modelo de referencia de ambos \textit{benchmarks} es una red recurrente con unidades GRU, y resulta llamativo lo difícil que ha sido superarla: los intentos de sustituirla por modelos de espacio de estados de tipo Mamba igualaron la pérdida CTC y la tasa de error de fonema, pero quedaron por detrás en tasa de error por palabra, y las variantes basadas en Transformer tampoco mejoraron de forma clara el modelo recurrente \cite{b2t24lessons}. Se ha sugerido que la tarea no depende de forma significativa de la información de largo alcance, lo que retiraría la principal ventaja de las arquitecturas atencionales. Trabajos posteriores han obtenido resultados competitivos con codificadores Conformer \cite{khanday2026}, de modo que la cuestión no está cerrada.

\subsection{Objetivo de entrenamiento y unidad lingüística}

El entrenamiento se realiza casi universalmente con la pérdida CTC \cite{graves2006ctc}, que permite aprender a partir de pares secuencia-etiqueta sin alineamiento temporal explícito, introduciendo un símbolo en blanco y marginalizando sobre todos los alineamientos compatibles con la etiqueta. Es la elección natural aquí, porque no se dispone de la correspondencia temporal entre actividad cortical y unidades emitidas: en un participante con parálisis no hay audio de referencia con el que alinear.

La unidad de salida dominante es el \textbf{fonema}, con la justificación neurocientífica de que la corteza motora del habla codifica articulación y no ortografía. Es la unidad del sistema de referencia de ambos \textit{benchmarks} y la que emplean casi todas las entradas, que además seleccionan modelo por tasa de error de fonema. Se han explorado alternativas de granularidad más fina: el equipo ganador del Brain-to-Text~'24 obtuvo su mejora principal empleando difonos, es decir, transiciones entre fonemas, con el argumento de que capturan mejor la coarticulación \cite{b2t24lessons}. En el extremo opuesto, algunos trabajos recientes han demostrado que la decodificación directa a nivel de \textbf{carácter} es viable de forma totalmente extremo a extremo \cite{khanday2026}, lo que resulta especialmente relevante para esta memoria.

\subsection{El modelo de lenguaje y la formación de palabras}

La salida del codificador es una secuencia de distribuciones sobre unidades sublexicales que, por sí sola, produce texto con una tasa de error inaceptable. La conversión a texto se realiza mediante una búsqueda que combina esa evidencia acústica con un modelo de lenguaje. El procedimiento estándar en los sistemas de referencia es un decodificador ponderado de estados finitos que integra un modelo $n$-grama de orden alto para generar una lista de hipótesis, seguido de una segunda etapa de reordenamiento con un modelo de lenguaje de gran tamaño \cite{b2t24lessons}.

Esta etapa aporta una parte muy sustancial del rendimiento final del sistema, hasta el punto de que la comparación entre codificadores sólo es significativa si se realiza con el mismo decodificador lingüístico. Tiene además un coste de recursos considerable: el modelo de 5-gramas empleado en el \textit{benchmark} requiere del orden de \SI{300}{\giga\byte} de memoria principal y el de 3-gramas unos \SI{60}{\giga\byte}, cifras fuera del alcance de un entorno de trabajo estándar, lo que ha llevado a varios participantes a entrenar modelos propios de orden reducido \cite{cheng2026}. Existen implementaciones más eficientes en memoria específicamente diseñadas para este escenario \cite{lightbeam}.

Conviene distinguir dos formas de emplear el modelo de lenguaje, porque la diferencia es central en esta memoria. En el \textbf{reordenamiento} (\textit{rescoring}), el modelo acústico genera una lista de hipótesis y el modelo de lenguaje se limita a ordenarla; su techo de rendimiento es, por construcción, la mejor hipótesis presente en la lista. En la \textbf{fusión superficial} (\textit{shallow fusion}), el modelo de lenguaje interviene durante la búsqueda en haz, guiando la expansión de hipótesis y pudiendo por tanto alcanzar transcripciones que nunca habrían aparecido en la lista $n$-mejor del modelo acústico. Comparar el rendimiento real con el límite \textit{oracle} de la lista $n$-mejor permite discriminar empíricamente cuál de los dos regímenes está operando.

\section{Los \textit{benchmarks} Brain-to-Text}
\label{sec:sota-benchmarks}

La existencia de conjuntos de datos públicos con particiones fijas y métrica común ha sido determinante para la aceleración reciente del campo.

El \textbf{Brain-to-Text~'24} se construyó sobre datos del participante T12: 12.100 frases registradas en 25 sesiones a lo largo de cuatro meses con matrices de 128 electrodos en la corteza motora ventral \cite{lightbeam}. Su modelo de referencia, la red recurrente con decodificación en dos etapas descrita anteriormente, alcanzó un \SI{9.7}{\percent} de tasa de error por palabra, mientras que la mejor entrada de la competición llegó al \SI{5.8}{\percent} \cite{b2t24lessons}.

El \textbf{Brain-to-Text~'25}, empleado en este trabajo, amplía considerablemente la escala: 10.948 frases del participante T15 registradas en 45 sesiones a lo largo de 20 meses con 256 microelectrodos, agrupados en cuatro matrices de alta densidad de 64 canales cada una, de las que se extraen 512 características por ventana de \SI{20}{\milli\second} \cite{b2t25,card2024}. El horizonte temporal de 20 meses convierte la no estacionariedad en un problema de primer orden, muy por encima de lo que planteaba el conjunto anterior.

Las entradas que encabezan la clasificación de este segundo \textit{benchmark} recurren de forma sistemática a conjuntos de modelos y a modelos de lenguaje de gran tamaño para la fusión de sus salidas: una de ellas combina 29 codificadores preentrenados inicializados con semillas distintas y emplea GPT-4 ajustado para fusionar las hipótesis de fonemas y frases de todos ellos \cite{crossspecies2025}. Este dato es relevante para interpretar cualquier resultado individual: la distancia entre un sistema de modelo único y la cabeza de la clasificación no es únicamente una diferencia de arquitectura, sino también de presupuesto de cómputo.

\section{Modelos de habla preentrenados}
\label{sec:sota-habla}

En paralelo, el reconocimiento automático del habla ha experimentado su propia transformación. El paradigma actual consiste en preentrenar un modelo de gran capacidad sobre volúmenes masivos de audio y adaptarlo después a la tarea concreta.

Dentro de este paradigma conviven dos familias. La primera es la del \textbf{aprendizaje autosupervisado}, representada por wav2vec~2.0, HuBERT y WavLM, que aprenden representaciones a partir de audio no etiquetado mediante objetivos contrastivos o de predicción enmascarada, y requieren una etapa posterior de ajuste supervisado. La segunda es la de la \textbf{supervisión débil a gran escala}, cuyo exponente es Whisper \cite{radford2023whisper}, entrenado sobre 680.000 horas de audio con transcripción de calidad heterogénea recogidas de la web, y que funciona directamente como sistema de reconocimiento multilingüe sin ajuste adicional.

El principal inconveniente de Whisper para la investigación es que ni sus datos de entrenamiento ni su procedimiento son plenamente reproducibles. \textbf{OWSM} (\textit{Open Whisper-style Speech Model}) \cite{peng2023owsm} nació precisamente para cubrir ese hueco: reproduce el estilo de entrenamiento de Whisper empleando exclusivamente conjuntos de datos públicos y una herramienta de código abierto, ESPnet \cite{watanabe2018espnet}, publicando datos, recetas y pesos. La versión v3.1 \cite{peng2024owsmv31} sustituye el codificador Transformer por la arquitectura E-Branchformer \cite{kim2023ebranchformer}, que combina una rama de autoatención con una rama convolucional para capturar simultáneamente dependencias globales y patrones locales, y mejora tanto la precisión como la velocidad de inferencia respecto a la versión anterior. Arquitectónicamente, OWSM es un modelo codificador-decodificador entrenado con un objetivo conjunto CTC-atención \cite{kim2017jointctc}, lo que resulta conveniente aquí: la rama CTC puede utilizarse de forma aislada, prescindiendo por completo del decodificador autorregresivo.

\section{Adaptación eficiente en parámetros}
\label{sec:sota-peft}

Ajustar la totalidad de los parámetros de un modelo preentrenado sobre un conjunto de datos pequeño es caro y propenso al sobreajuste. Las técnicas de adaptación eficiente en parámetros (\textit{Parameter-Efficient Fine-Tuning}, PEFT) abordan el problema modificando únicamente una fracción reducida de los pesos.

La más extendida es \textbf{LoRA} \cite{hu2022lora}, que congela los pesos originales e introduce, en paralelo a las matrices de proyección seleccionadas, un par de matrices de rango reducido cuyo producto se suma a la salida. La hipótesis subyacente es que la actualización necesaria para adaptarse a una tarea nueva tiene rango intrínseco bajo. Alternativas habituales son los módulos adaptadores insertados entre capas, el ajuste de prefijos y, en su versión más simple, el \textbf{descongelado parcial}: mantener congeladas las capas inferiores y ajustar únicamente las superiores.

Es importante subrayar un punto que condiciona el diseño experimental de esta memoria: \textbf{LoRA sólo tiene sentido sobre un modelo base congelado}. Aplicado sobre un codificador completamente descongelado, añade parámetros entrenables sin aportar ninguna restricción, por lo que no cabe esperar de él ningún beneficio. Cualquier comparación entre estrategias de adaptación debe controlar esta interacción explícitamente.

\section{Transferencia de modelos de habla a señal neuronal}
\label{sec:sota-transferencia}

La idea de reutilizar un modelo entrenado sobre audio para decodificar actividad cerebral no es nueva, y conviene revisar con detalle los trabajos previos porque delimitan con precisión qué queda por responder.

Una primera línea emplea modelos de habla como \textbf{espacio de representación objetivo}. Kim \textit{et al.} \cite{braintalker} entrenan un codificador que proyecta señal ECoG al espacio latente de wav2vec~2.0, para sintetizar después voz inteligible en un escenario de muy pocos datos. En una variante próxima, se ha reentrenado la arquitectura de wav2vec sobre registros ECoG no etiquetados y se han empleado las representaciones resultantes como entrada de un decodificador supervisado, con mejoras consistentes respecto a usar la señal original y con transferencia positiva incluso entre pacientes distintos \cite{ecogssl2024}. En estos trabajos, sin embargo, lo que se transfiere es la \textit{arquitectura} o el \textit{objetivo autosupervisado}, no los pesos aprendidos sobre audio.

La línea directamente relevante para esta memoria es la que \textbf{transfiere los pesos preentrenados sobre habla} sustituyendo la entrada acústica por señal neuronal. Fiedler \textit{et al.} \cite{fiedler2025wav2vec2brain} reemplazan el extractor de características de audio de wav2vec~2.0 por un extractor de características neuronales entrenado desde cero y comparan sistemáticamente tres condiciones sobre los datos del participante T12: ajuste completo partiendo de los pesos preentrenados, entrenamiento desde cero con la misma arquitectura, y modelo preentrenado congelado con únicamente el extractor entrenable. El ajuste completo con pesos preentrenados obtiene la mejor tasa de error por carácter, \SI{18.54}{\percent}, superando en \num{20.46} puntos porcentuales al entrenamiento desde cero y en \num{15.92} puntos a la variante congelada. Es, hasta donde alcanza esta revisión, la evidencia más directa de que existe transferencia real desde el reconocimiento del habla hacia la decodificación neuronal.

Más recientemente, BrainWhisperer \cite{brainwhisperer} lleva esta idea a Whisper y al escenario intracortical de 256 canales. Su punto de partida es un hallazgo de interpretabilidad: el codificador de Whisper aprende representaciones selectivas de fonema con atención localizada, lo que sugiere que sus capas iniciales son un extractor articulatorio reutilizable. Sobre esa base construyen un modelo con pérdida híbrida, CTC de fonemas tomada de la tercera capa del codificador y entropía cruzada sobre unidades de palabra, atención propia con ventana para capturar continuidad articulatoria, proyecciones jerárquicas de bajo rango por mes y por día frente a la no estacionariedad, e incrustaciones específicas de sujeto que habilitan el entrenamiento conjunto entre participantes. El sistema se sitúa entre las mejores entradas del Brain-to-Text~'25.
% TODO: verifica en el propio artículo el puesto exacto y el WER (la cifra de
% 1,78 % y el 4.º puesto de 466 la he visto sólo en prensa, no en el paper).

Existe además una tercera línea, complementaria, que preentrena \textbf{modelos fundacionales específicos de señal neuronal} en lugar de reutilizar modelos de habla: encoders entrenados sobre múltiples tareas motoras, participantes e incluso especies, que se ajustan después a la tarea de habla \cite{crossspecies2025}. Esta aproximación evita el desajuste de modalidad, a costa de requerir la agregación de grandes volúmenes de datos neuronales.

\section{Posicionamiento de este trabajo}
\label{sec:sota-posicionamiento}

De la revisión anterior se desprenden tres observaciones que delimitan la aportación de esta memoria.

En primer lugar, \textbf{la transferencia desde modelos de habla a señal neuronal está demostrada, pero poco caracterizada}. El trabajo de Fiedler \textit{et al.} \cite{fiedler2025wav2vec2brain} establece que existe, sobre wav2vec~2.0 y el participante T12. BrainWhisperer \cite{brainwhisperer} la explota con éxito sobre Whisper, pero como parte de un sistema con múltiples modificaciones simultáneas orientado a maximizar el rendimiento en la competición, lo que hace difícil atribuir la mejora a un componente concreto. Falta, entre ambos, un estudio de ablación con presupuesto de cómputo igualado que aísle cada decisión de diseño por separado sobre el conjunto de datos más reciente.

En segundo lugar, \textbf{la elección de la unidad lingüística de salida rara vez se justifica empíricamente}. El fonema se adopta por argumento neurocientífico y por continuidad con el sistema de referencia, mientras que el carácter aparece en trabajos aislados \cite{khanday2026}. No se ha encontrado una comparación sistemática de granularidades sobre un mismo modelo, un mismo presupuesto y un mismo eje de evaluación, que es la única forma de contrastar ese argumento en lugar de asumirlo.

En tercer lugar, \textbf{los resultados publicados son difíciles de reproducir con recursos modestos}. Los sistemas de cabecera combinan decenas de modelos y modelos de lenguaje de gran tamaño, y hasta el modelo $n$-grama de referencia excede la memoria disponible en un entorno estándar \cite{cheng2026}. Documentar qué parte del rendimiento sigue siendo accesible con una única GPU de gama media y un modelo de lenguaje entrenado sobre el propio corpus tiene, por tanto, valor metodológico.

Este trabajo se sitúa en la intersección de las tres observaciones. Adapta OWSM v3.1 \cite{peng2024owsmv31} —y no Whisper, por ser un modelo íntegramente reproducible, con datos y recetas públicos— a los datos del Brain-to-Text~'25, y en lugar de perseguir la mejor cifra posible, ejecuta un conjunto de ablaciones controladas, con una única variable modificada por experimento y presupuesto de cómputo idéntico, orientadas a responder las tres preguntas de investigación formuladas en el \cref{cap:introduccion}: qué transfiere del preentrenamiento, a qué granularidad conviene leer la corteza, y dónde reside el conocimiento lingüístico del sistema.

Como punto de referencia directo para los resultados obtenidos, resulta especialmente útil el trabajo de Khanday \textit{et al.} \cite{khanday2026}, que emplea el mismo conjunto de datos, la misma unidad de salida a nivel de carácter, una capa de alineamiento específica de sesión análoga a la utilizada aquí y decodificación CTC voraz sin modelo de lenguaje externo, alcanzando una tasa de error por carácter del \SI{23.80}{\percent} sobre la partición de validación. Al compartir métrica, datos y condiciones de decodificación, constituye la comparación más informativa disponible para un codificador entrenado desde cero, y por tanto una referencia natural frente a la que contrastar la hipótesis de transferencia que motiva esta memoria.

% =====================================================================
%  CAPÍTULO 3 — CONJUNTO DE DATOS
% =====================================================================
% Requiere en el preámbulo: booktabs, enumitem, cleveref, siunitx, multirow
% =====================================================================

\chapter{Conjunto de datos}
\label{cap:datos}

La elección del conjunto de datos condiciona el resto del trabajo: determina qué preguntas pueden formularse, qué arquitecturas son viables y con qué resultados publicados es posible compararse. Este capítulo justifica esa elección frente a las alternativas evaluadas, describe en detalle el conjunto finalmente seleccionado y analiza los retos específicos que plantea, varios de los cuales explican decisiones de diseño posteriores.

\section{Criterios de selección}
\label{sec:criterios-datos}

Antes de fijar el conjunto de trabajo se evaluaron cuatro corpus públicos de decodificación de habla a partir de señal cerebral. Los criterios aplicados fueron los siguientes:

\begin{enumerate}[label=\textbf{C\arabic*.}, leftmargin=*, itemsep=0.3em]
  \item \textbf{Tarea compatible con el modelo.} El sistema propuesto reutiliza un codificador de reconocimiento del habla entrenado sobre secuencias largas. Requiere, por tanto, un corpus de \textbf{frases continuas}, no de palabras o sílabas aisladas: un modelo diseñado para modelar dependencias temporales de varios segundos no puede evaluarse sobre estímulos de dos segundos.
  \item \textbf{Salida textual.} El objetivo es \textit{brain-to-text}, no \textit{brain-to-speech}. Los corpus orientados a la reconstrucción de audio requieren una cadena de síntesis (mel-espectrograma más vocoder) y un tipo de evaluación distinto, ajenos a este trabajo.
  \item \textbf{Volumen suficiente.} El ajuste de un modelo de más de cien millones de parámetros exige un número de muestras que descarta los corpus de pocos minutos por sujeto.
  \item \textbf{Relación señal-ruido.} La hipótesis de partida —que un codificador de habla puede extraer estructura temporal de la señal neuronal— sólo es contrastable si la señal contiene esa estructura. Las modalidades de baja resolución espacial la comprometen.
  \item \textbf{Comparabilidad.} La existencia de particiones fijas, métrica oficial y resultados publicados permite situar los resultados propios en un contexto objetivo, en lugar de reportarlos de forma aislada.
  \item \textbf{Coste de acceso y preparación.} Formato documentado, código de carga disponible y datos ya organizados en particiones reducen el tiempo dedicado a ingeniería de datos, escaso en un trabajo de esta duración.
\end{enumerate}

\section{Alternativas evaluadas}
\label{sec:alternativas-datos}

La \cref{tab:datasets-comparativa} resume los cuatro conjuntos considerados.

\begin{table}[htbp]
  \centering
  \caption{Comparativa de los conjuntos de datos evaluados: señal y tarea.}
  \label{tab:datasets-comparativa}
  \small
  \begin{tabular}{p{2.6cm} p{4.4cm} p{2.0cm} p{3.2cm}}
    \toprule
    \textbf{Conjunto} & \textbf{Tipo de señal} & \textbf{Sujetos} & \textbf{Nivel de tarea} \\
    \midrule
    UGR-MINDVOICE   & EEG (64 canales)                     & 15 & Palabras y sílabas \\
    sEEG2Speech     & sEEG                                 & 3  & Frases (5--7 palabras) \\
    SingleWord (NL) & sEEG (1.103 canales)                 & 10 & Palabras aisladas \\
    Brain-to-Text '25 & Microelectrodos intracorticales (256) & 1 & Frases continuas \\
    \bottomrule
  \end{tabular}

  \vspace{1em}

  \caption{Comparativa de los conjuntos de datos evaluados: volumen y formato.}
  \label{tab:datasets-comparativa2}
  \small
  \begin{tabular}{p{2.6cm} p{3.6cm} p{2.6cm} p{3.4cm}}
    \toprule
    \textbf{Conjunto} & \textbf{Volumen} & \textbf{Salida objetivo} & \textbf{Formato} \\
    \midrule
    UGR-MINDVOICE   & 60 palabras, 30 pseudo\-palabras, 95 sílabas & Clasificación o texto & MNE/EDF \\
    sEEG2Speech     & 100 frases, 10--20 min por sujeto & Audio & Propio (OSF) \\
    SingleWord (NL) & 100 palabras, $\sim$5 min por sujeto & Espectrograma o audio & BIDS / NWB \\
    Brain-to-Text '25 & 10.948 frases, 45 sesiones en 20 meses & Texto & HDF5 con particiones \\
    \bottomrule
  \end{tabular}
\end{table}

\subsection{UGR-MINDVOICE}

Corpus multimodal de EEG de superficie con 15 participantes nativos de español, que incluye tanto habla pronunciada (\textit{overt}) como imaginada (\textit{covert}), con estímulos de 95 sílabas consonante-vocal que cubren el inventario fonético del español, 60 palabras de seis categorías semánticas y 30 pseudopalabras. Es el único corpus en español para esta tarea y su diseño experimental es cuidadoso, con datos y código públicos.

Fue descartado por dos razones. La primera es el nivel de la tarea: los estímulos son palabras y sílabas aisladas, incompatibles con el criterio C1. La segunda es la modalidad: el EEG de superficie tiene una resolución espacial y una relación señal-ruido muy inferiores a las de los registros invasivos, lo que compromete el criterio C4 y, con él, la \textbf{interpretabilidad de un resultado negativo}. Si el sistema fracasase, no sería posible distinguir entre ausencia de transferencia desde el modelo de habla y ausencia de señal decodificable en la entrada, que son conclusiones muy distintas. El descarte responde, por tanto, al diseño experimental de este trabajo concreto y no constituye un juicio sobre la calidad del corpus, cuyo diseño es cuidadoso y cuya orientación a un idioma distinto del inglés cubre un hueco real.

\subsection{sEEG2Speech}

Corpus de estereoelectroencefalografía en tres pacientes con epilepsia, que leen 100 frases de cinco a siete palabras en neerlandés. La señal se filtra en la banda gamma alta (70--\SI{170}{\hertz}) y el objetivo es la reconstrucción directa de audio mediante mel-espectrogramas y un vocoder.

Cumple el criterio C1 pero incumple el C2 —el objetivo es audio, no texto— y sobre todo el C3: entre diez y veinte minutos por sujeto, de los cuales una fracción considerable es silencio, es un volumen incompatible con el ajuste de un modelo de gran tamaño.

\subsection{SingleWordProductionDutch}

Diez pacientes con epilepsia y 1.103 canales sEEG en total, leyendo 100 palabras aisladas en neerlandés, con formato BIDS bien documentado y una cadena de referencia basada en filtrado en gamma alta, reducción dimensional y regresión hacia el log-mel espectrograma.

Es el corpus con más participantes de los evaluados, lo que lo haría interesante para estudios de generalización entre sujetos, pero incumple C1 (palabras aisladas), C2 (salida de audio) y C3 (unos cinco minutos por sujeto).

\subsection{Brain-to-Text '25}

Es el único de los cuatro que satisface todos los criterios simultáneamente: frases continuas, salida textual, volumen de casi once mil frases, registros intracorticales de alta relación señal-ruido, particiones oficiales con clasificación pública y resultados publicados, y datos organizados en HDF5 con código de carga proporcionado por los autores. Es, en consecuencia, el conjunto empleado en este trabajo.

Conviene señalar que esta elección tiene un coste evidente: \textbf{un único participante}. Las conclusiones obtenidas son específicas de T15, de la localización de sus implantes y del inglés, y no pueden extrapolarse sin verificación a otros sujetos. Se asume este compromiso porque los tres corpus alternativos, aunque más diversos en participantes, no permiten formular la pregunta de investigación de este trabajo.

\section{El conjunto Brain-to-Text '25}
\label{sec:b2t25}

\subsection{Origen y participante}

Los datos proceden del ensayo clínico descrito por Card \textit{et al.} \cite{card2024} y se publicaron posteriormente como \textit{benchmark} abierto y competición \cite{b2t25}. El participante, identificado como T15, es un varón con esclerosis lateral amiotrófica avanzada y disartria severa, con el habla ininteligible pero con la corteza motora del habla funcionalmente intacta.

El protocolo es una tarea de copia: se presenta una frase en pantalla y, tras una señal de inicio, el participante intenta pronunciarla a su propio ritmo. Se registra la actividad neuronal durante ese intento. Es importante subrayar que \textbf{no existe audio de referencia} con el que alinear la señal, ya que la parálisis impide la producción de habla inteligible; de ahí que el entrenamiento deba recurrir necesariamente a un objetivo que no requiera alineamiento explícito, como CTC.

\subsection{Adquisición y características neuronales}

El registro se realiza mediante cuatro matrices de microelectrodos de alta densidad de 64 canales cada una, implantadas en la corteza motora del habla del hemisferio izquierdo, lo que suma 256 electrodos \cite{b2t25}.
% TODO: consulta Card et al. 2024 para citar las cuatro áreas concretas
% (giro precentral izquierdo; se mencionan 6v ventral y dorsal, área 4 y 55b
% en la literatura secundaria, pero verifícalo en la fuente primaria).

De cada electrodo se extraen dos características, lo que produce un vector de \textbf{512 dimensiones por ventana temporal}:

\begin{description}
  \item[Cruces de umbral (\textit{threshold crossings}).] Recuento de eventos en que la señal filtrada cruza un umbral fijado en $-4{,}5$ veces el valor eficaz del ruido del canal. Es una aproximación a la tasa de disparo de la población neuronal próxima al electrodo, sin necesidad de clasificar espigas individuales.
  \item[Potencia en la banda de espigas (\textit{spike band power}).] Energía de la señal en la banda de frecuencia asociada a la actividad de espiga, que aporta información complementaria y es más robusta frente a variaciones del umbral.
\end{description}

Ambas se agregan en ventanas de \SI{20}{\milli\second}, produciendo una secuencia a \SI{50}{\hertz}, y se normalizan mediante puntuación z dentro de cada bloque de registro \cite{lightbeam}. Esta normalización por bloque es la primera línea de defensa frente a la no estacionariedad, pero, como se verá, resulta insuficiente por sí sola.

\subsection{Organización de los datos}

Los datos se distribuyen en ficheros HDF5 organizados por sesión de registro. Cada sesión se identifica por su fecha (\texttt{t15.AAAA.MM.DD}) y contiene hasta tres ficheros, correspondientes a las particiones de entrenamiento, validación y prueba. No todas las sesiones aportan a las tres particiones: algunas contienen únicamente datos de entrenamiento.

Dentro de cada fichero, cada ensayo se almacena como un grupo independiente con los siguientes elementos \cite{b2t25}:

\begin{description}
  \item[\texttt{input\_features}.] Matriz de características neuronales de dimensiones $T \times 512$, donde $T$ es el número de ventanas del ensayo.
  \item[\texttt{seq\_class\_ids}.] Secuencia de identificadores de fonema correspondiente a la frase, es decir, la etiqueta fonémica de referencia.
  \item[\texttt{transcription}.] Transcripción de la frase.
  \item[Atributos.] Número de ventanas temporales, longitud de la secuencia de etiquetas, etiqueta textual de la frase, e identificadores de sesión, bloque y ensayo.
\end{description}

Es relevante para el diseño experimental que el conjunto \textbf{ya incluye la secuencia de fonemas de referencia}: el experimento de decodificación a nivel de fonema no requiere ningún sistema de conversión grafema-fonema ni diccionario de pronunciación adicional, lo que reduce sustancialmente su coste de preparación.

\subsection{Particiones y composición por corpus}

La \cref{tab:b2t25-splits} resume la composición del conjunto. El reparto es de 8.072 ensayos de entrenamiento, 1.426 de validación y 1.450 de prueba, estos últimos sin etiquetas, reservados para la evaluación oficial \cite{lightbeam}. Los ensayos se agrupan en bloques de entre 5 y 50 frases; el conjunto de entrenamiento comprende 198 bloques repartidos por las 45 sesiones, mientras que validación y prueba suman 67 bloques procedentes de 41 de esas 45 sesiones.

\begin{table}[htbp]
  \centering
  \caption{Composición del conjunto Brain-to-Text '25 por corpus de origen. Datos tomados de \cite{lightbeam}.}
  \label{tab:b2t25-splits}
  \small
  \begin{tabular}{lrrl}
    \toprule
    \textbf{Corpus de origen} & \textbf{Bloques} & \textbf{Frases} & \textbf{Partición} \\
    \midrule
    Switchboard              & 182 & 7.672 & Entrenamiento y validación/prueba \\
    Palabras frecuentes      &  48 & 2.050 & Entrenamiento y validación/prueba \\
    Vocabulario de 50 palabras & 25 & 1.091 & Solo entrenamiento \\
    OpenWebText              &   5 &   172 & Solo validación/prueba \\
    Frases Harvard           &   3 &    98 & Solo validación/prueba \\
    Secuencias aleatorias    &   2 &    79 & Solo validación/prueba \\
    \bottomrule
  \end{tabular}
\end{table}

Esta tabla merece atención porque revela un rasgo del \textit{benchmark} que no es evidente a primera vista y que tiene consecuencias directas sobre la interpretación de los resultados. \textbf{Tres de los seis corpus de origen aparecen exclusivamente en validación y prueba, y nunca en entrenamiento}: OpenWebText, las frases Harvard y las secuencias aleatorias de palabras suman 349 frases que el modelo evalúa sin haber visto material lingüístico de esa procedencia. En sentido inverso, el vocabulario cerrado de 50 palabras aparece únicamente en entrenamiento.

Es decir, la partición de validación no es una muestra homogénea de la de entrenamiento: incorpora por diseño un \textbf{desplazamiento de dominio lingüístico}. Las secuencias aleatorias de palabras son, además, el caso adverso extremo para cualquier modelo de lenguaje, puesto que carecen por construcción de estructura sintáctica explotable. Esto explica parcialmente la marcada divergencia entre pérdida de entrenamiento y pérdida de validación que reportan los participantes de la competición, y obliga a matizar cualquier atribución de esa divergencia únicamente al sobreajuste del codificador.

\subsection{Estrategias de habla no etiquetadas}

Un último rasgo relevante: la mayor parte de los datos corresponde a intentos de habla \textit{vocalizada}, pero una fracción corresponde a intentos de habla \textit{silente}, y los organizadores decidieron deliberadamente \textbf{no proporcionar la etiqueta de estrategia} de cada bloque, con el objetivo de aproximar mejor las condiciones reales de uso \cite{b2t25}.

Se trata, por tanto, de una variable oculta que modula la señal y que el modelo debe absorber implícitamente. Ninguna de las variantes evaluadas en este trabajo la modela de forma explícita, lo que constituye una limitación reconocida y, a la vez, una línea de mejora identificable.

\section{Análisis exploratorio}
\label{sec:eda}

% TODO (Kiril): esta sección la tienes que rellenar con tus propios números y
% figuras a partir del notebook de análisis. Te dejo la estructura y qué
% debería mostrar cada figura, porque es la sección que da solidez a todo lo
% que viene después y es lo primero que va a mirar un tribunal técnico.

\subsection{Distribución de longitudes}

% Figura sugerida: histograma de T (nº de ventanas por ensayo) y de la
% longitud de la transcripción en caracteres y en fonemas.
% Comenta: rango, mediana, cola larga, e implicación para el batching y el
% coste de memoria (los ensayos largos dominan el consumo).

\subsection{Variabilidad entre sesiones}

% Figura sugerida (la más importante del capítulo): estadísticos de las
% características neuronales (media y desviación por canal) a lo largo de las
% 45 sesiones, ordenadas cronológicamente. Si hay deriva, se verá.
% Alternativa complementaria: matriz de distancias entre sesiones sobre la
% distribución de features. Esta figura es la que justifica la capa de sesión.

\subsection{Cobertura léxica y ejemplos por clase}

% Figura sugerida, y es de coste casi cero: número de ejemplos de
% entrenamiento por clase para cada unidad de salida candidata (carácter,
% fonema, subpalabra, palabra). En escala logarítmica.
% Esta figura sola justifica todo el bloque experimental de granularidad:
% con ~8.000 frases, el carácter tiene miles de ejemplos por clase mientras
% que buena parte del vocabulario de palabras aparece una o dos veces.

\subsection{Distribución temporal del registro}

% Figura sugerida: número de ensayos por sesión a lo largo de los 20 meses,
% distinguiendo partición. Sirve para mostrar que los datos no están
% uniformemente repartidos en el tiempo.

\section{Retos que plantea el conjunto y consecuencias de diseño}
\label{sec:retos-datos}

Del análisis anterior se derivan cinco retos que condicionan directamente las decisiones tomadas en los capítulos siguientes:

\begin{enumerate}[label=\textbf{R\arabic*.}, leftmargin=*, itemsep=0.4em]
  \item \textbf{No estacionariedad a lo largo de 20 meses.} Es el horizonte temporal más amplio de los \textit{benchmarks} disponibles y multiplica por cinco el del conjunto anterior. La normalización por bloque no basta: motiva la incorporación de una transformación afín específica de cada sesión, y convierte la evaluación sobre sesiones no vistas en un experimento con valor propio.
  \item \textbf{Ausencia de alineamiento temporal.} Al no existir audio de referencia, no se dispone de correspondencia entre ventanas de señal y unidades emitidas. Determina el uso de CTC como objetivo de entrenamiento.
  \item \textbf{Escala de datos frente a escala de modelo.} Poco más de ocho mil frases de entrenamiento para un modelo de más de cien millones de parámetros sitúan el problema en un régimen de alto riesgo de sobreajuste. Un corolario aparente de este hecho es que las unidades de salida con vocabularios grandes deberían resultar inviables, puesto que repartir ese número de frases entre decenas de miles de clases de subpalabra deja a la mayoría con un puñado de ejemplos o ninguno. Conviene anticipar que \textbf{el experimento desmiente ese corolario} (\cref{sec:resultados-B}): el reparto de ejemplos por clase no es el factor determinante, porque las clases no observadas no penalizan bajo CTC y la longitud de la secuencia objetivo pesa más que la densidad del vocabulario.
  \item \textbf{Desplazamiento de dominio lingüístico entre particiones.} Los corpus presentes solo en validación penalizan a los sistemas que dependan en exceso del modelo de lenguaje, y en particular las secuencias aleatorias de palabras actúan como caso adverso deliberado. Debe tenerse en cuenta al interpretar la diferencia entre las métricas de entrenamiento y validación.
  \item \textbf{Variable oculta de estrategia de habla.} La mezcla no etiquetada de intento vocalizado y silente introduce una fuente de variabilidad no modelada que se suma a la variabilidad entre sesiones.
\end{enumerate}

Estos cinco puntos, junto con el marco de referencia establecido en el \cref{cap:estado-del-arte}, definen el espacio de diseño del sistema que se describe a continuación.

\include{capitulos/04_metodologia}
\include{capitulos/05_experimentos}
\include{capitulos/06_resultados}
\include{capitulos/07_conclusiones}

\appendix
\include{capitulos/A_anexos}

\printbibliography[heading=bibintoc,title={Bibliografía}]

\end{document}


% ---------------------------------------------------------------------
% (4) TRUCO: compilar solo el capítulo en el que estás trabajando
%     Con \includeonly, LaTeX salta el resto pero mantiene numeración
%     de páginas, referencias cruzadas y citas correctas.
%     Lo pones en el PREÁMBULO y lo comentas cuando quieras el PDF entero:
% ---------------------------------------------------------------------
\includeonly{capitulos/01_introduccion}


% ---------------------------------------------------------------------
% (5) ETIQUETAS QUE TIENES QUE PONER EN CADA CAPÍTULO
%     El capítulo 1 las referencia con \cref{...}. Si no existen,
%     te saldrán interrogantes en el PDF.
% ---------------------------------------------------------------------
% \chapter{Estado del arte}      \label{cap:estado-del-arte}
% \chapter{Conjunto de datos}    \label{cap:datos}
% \chapter{Metodología}          \label{cap:metodologia}
% \chapter{Diseño experimental}  \label{cap:experimentos}
% \chapter{Resultados}           \label{cap:resultados}
% \chapter{Conclusiones}         \label{cap:conclusiones}
